\documentclass[11pt]{article}

% Answers and hints for the practice problems of Lecture 2
\usepackage[margin=1in]{geometry}
\usepackage[T1]{fontenc}
\usepackage{lmodern}
\usepackage{amsmath,amssymb,amsthm}
\usepackage{enumitem}
\usepackage{xr-hyper}
\usepackage{hyperref}
% Problem numbers are read from practice2.aux, so compile practice2.tex first.
\externaldocument{practice2}[practice2.pdf]
\usepackage{xcolor}

\title{Lecture 2: Practice Problems --- Answers and Hints}
\author{Linear Algebra}
\date{\today}

\newcommand{\norm}[1]{\lVert #1 \rVert}
\newcommand{\ip}[2]{\langle #1, #2 \rangle}

\DeclareMathOperator{\im}{im}
\DeclareMathOperator{\rank}{rank}
\DeclareMathOperator{\Null}{null}
\DeclareMathOperator{\col}{col}
\DeclareMathOperator{\row}{row}
\DeclareMathOperator{\Span}{span}
\DeclareMathOperator{\tr}{tr}
\DeclareMathOperator{\diag}{diag}

\begin{document}

\maketitle

\section*{Norms}

\begin{description}[leftmargin=0pt, itemsep=0.3em]
    \item[\ref{pr:norms-compute}.] (a) $9$, $5$, $4$. (b) $7$, $5$, $337^{1/4} \approx 4.28$, $4$: the $p$-norms decrease towards the $\ell_\infty$ norm. (c) $\max_i |x_i|^p \leq \sum_i |x_i|^p \leq n \max_i |x_i|^p$; take $p$-th roots, and use $n^{1/p} \to 1$.
    \item[\ref{pr:which-norms}.] (a) A norm; its unit ball is the diamond with vertices $(\pm 1, 0)$ and $(0, \pm\frac{1}{2})$. (b) Not absolutely homogeneous: $f(2\mathbf{e}_2) = 4 \neq 2f(\mathbf{e}_2)$. (c) Not positive definite: $f(1, 1) = 0$. (d) A norm (a sum of two functions that satisfy the triangle inequality, and which is $0$ only for $\mathbf{x} = \mathbf{0}$). (e) The triangle inequality fails: $f(\mathbf{e}_1 + \mathbf{e}_2) = 4 > 2 = f(\mathbf{e}_1) + f(\mathbf{e}_2)$.
\end{description}

\section*{Inner Products}

\begin{description}[leftmargin=0pt, itemsep=0.3em]
    \item[\ref{pr:which-ip}.] (a) Yes, $A = \diag(1, 2)$. (b) No: $\ip{\mathbf{e}_2}{\mathbf{e}_2} = -1$. (c) Yes, $A = \begin{pmatrix} 2 & 1 \\ 1 & 1 \end{pmatrix}$ ($2 > 0$, $\det = 1 > 0$). (d) No: $\ip{\mathbf{e}_1}{\mathbf{e}_1} = 0$. (e) No: not symmetric. (f) No: $\ip{\mathbf{x}}{\mathbf{x}} = 0$ for $\mathbf{x} = (1, -2)^\top$, since $\det = 0$.
    \item[\ref{pr:ip-gram}.] (a) $\ip{\mathbf{e}_1}{\mathbf{e}_2} = 1$, $\norm{\mathbf{e}_1} = \sqrt{2}$, $\norm{\mathbf{e}_2} = 1$, angle $45^\circ$. (b) $(1, -2)^\top$. (c) $\frac{1}{\sqrt{2}}(1, 0)^\top$ and $\big(-\frac{1}{\sqrt{2}}, \sqrt{2}\big)^\top$.
    \item[\ref{pr:spd-param}.] (a) $|t| < 2$; $t > 1$. (b) The upper left determinants are $2$, $3$, $4$.
    \item[\ref{pr:frobenius-ip}.] (a) $\tr(SK) = \tr\big((SK)^\top\big) = \tr(K^\top S^\top) = -\tr(KS) = -\tr(SK)$, so $\tr(S^\top K) = \tr(SK) = 0$. (b) Pythagoras for the Frobenius inner product.
\end{description}

\section*{Lengths and Distances}

\begin{description}[leftmargin=0pt, itemsep=0.3em]
    \item[\ref{pr:cs-apps}.] (a) $2x + 3y + 6z \leq \norm{(2, 3, 6)} \cdot 1 = 7$, attained at $\frac{1}{7}(2, 3, 6)$. (b) Apply Cauchy--Schwarz to $(\sqrt{a}, \sqrt{b}, \sqrt{c})$ and $(\frac{1}{\sqrt{a}}, \frac{1}{\sqrt{b}}, \frac{1}{\sqrt{c}})$; equality for $a = b = c$. (c) $\norm{(1, 2, 2)}^2 = 9$.
    \item[\ref{pr:polarization}.] (a) $\frac{1}{4}(49 - 25) = 6$. (b) By the parallelogram law, $49 + 25 = 2 \cdot 25 + 2\norm{y}^2$, so $\norm{y} = 2\sqrt{3}$.
    \item[\ref{pr:dist}.] (a) The difference is $(-3, 2, 6)$: $11$, $7$, $6$, and $\sqrt{9 + 16 + 36} = \sqrt{61}$. (b) Apply the triangle inequality to $x = (x - y) + y$ and to $y = (y - x) + x$.
\end{description}

\section*{Angles and Orthogonality}

\begin{description}[leftmargin=0pt, itemsep=0.3em]
    \item[\ref{pr:angles}.] (a) $90^\circ$. (b) $\cos\omega = \frac{4}{9}$, $\omega \approx 63.6^\circ$. (c) $\cos\omega = \frac{1}{\sqrt{3}}$, $\omega \approx 54.7^\circ$. (d) $(1, 1, 0)$ and $(1, 0, 1)$: $60^\circ$.
    \item[\ref{pr:orth-to-two}.] $\Span\{(-1, 1, 1)^\top\}$.
    \item[\ref{pr:householder-q}.] (a) The columns are orthonormal, and $Q^2 = Q^\top Q = I_3$. (b) $Q\mathbf{x} = \frac{1}{3}(5, 1, 1)^\top$, with length $\sqrt{3} = \norm{\mathbf{x}}$. (c) $Q$ is the reflection in the plane orthogonal to $\mathbf{u}$, i.e., the plane $\mathbf{x} - y - z = 0$; $\det(Q) = -1$.
    \item[\ref{pr:cosine-docs}.] (a) $1$ and $\frac{1}{\sqrt{60}} \approx 0.13$. (b) $\sqrt{24} \approx 4.90$ and $\sqrt{14} \approx 3.74$. (c) Cosine similarity says $d_2$, the Euclidean distance says $d_3$. Since $d_2 = 3d_1$ is the same document repeated three times, the cosine similarity gives the more sensible answer.
\end{description}

\section*{Orthonormal Bases}

\begin{description}[leftmargin=0pt, itemsep=0.3em]
    \item[\ref{pr:gs-haar}.] (a) $\frac{1}{2}(1, 1, 1, 1)$, $\frac{1}{2}(1, 1, -1, -1)$, $\frac{1}{\sqrt{2}}(1, -1, 0, 0)$. (b) $\frac{1}{\sqrt{2}}(0, 0, 1, -1)$.
    \item[\ref{pr:onb-coords}.] (b) $\ip{\mathbf{v}}{\mathbf{b}_1} = 3$, $\ip{\mathbf{v}}{\mathbf{b}_2} = 0$, $\ip{\mathbf{v}}{\mathbf{b}_3} = 3$, so $\mathbf{v} = 3\mathbf{b}_1 + 3\mathbf{b}_3$; and $3^2 + 0^2 + 3^2 = 18 = \norm{\mathbf{v}}^2$.
    \item[\ref{pr:qr}.] $Q$ has the columns $\frac{1}{\sqrt{2}}(1, 1, 0)^\top$ and $\frac{1}{\sqrt{6}}(1, -1, 2)^\top$, and $R = \begin{pmatrix} \sqrt{2} & \frac{1}{\sqrt{2}} \\ 0 & \frac{\sqrt{6}}{2} \end{pmatrix}$.
\end{description}

\section*{Orthogonal Complements}

\begin{description}[leftmargin=0pt, itemsep=0.3em]
    \item[\ref{pr:complement}.] $U^\perp$ is the null space of the matrix with rows $(1, 2, 0, 1)$ and $(0, 1, 1, 0)$; a basis is $(2, -1, 1, 0)^\top$, $(-1, 0, 0, 1)^\top$.
    \item[\ref{pr:four-subspaces}.] (a) $\col(A)$ is spanned by $(1, 2, 3)^\top$; $\Null(A^\top)$ by $(-2, 1, 0)^\top$ and $(-3, 0, 1)^\top$; $\row(A)$ by $(1, 2)^\top$; and $\Null(A)$ by $(-2, 1)^\top$. (b) Exactly when $\mathbf{b}$ is orthogonal to $\Null(A^\top)$, i.e., when $\mathbf{b}$ is a multiple of $(1, 2, 3)^\top$.
    \item[\ref{pr:classifier}.] (b) $\mathbf{w}^\top \mathbf{x} + b$ is $2$, $-1$, $-1$, so the classes are $+1$, $-1$, $-1$, and the distances are $\frac{2}{\sqrt{5}}$, $\frac{1}{\sqrt{5}}$, $\frac{1}{\sqrt{5}}$. The point $(1, 1)$ has the largest margin.
    \item[\ref{pr:hyperplane-form}.] $\frac{x}{1} + \frac{y}{2} + \frac{z}{3} = 1$, i.e., $6x + 3y + 2z = 6$, with normal vector $(6, 3, 2)$ of length $7$; the distance is $\frac{6}{7}$.
\end{description}

\section*{Inner Products of Functions}

\begin{description}[leftmargin=0pt, itemsep=0.3em]
    \item[\ref{pr:func-ip}.] (a) $\ip{1}{x} = \frac{1}{2}$, $\norm{1} = 1$, $\norm{x} = \frac{1}{\sqrt{3}}$, so $\cos\omega = \frac{\sqrt{3}}{2}$ and $\omega = 30^\circ$. (b) $1$ and $x - \frac{1}{2}$. (c) $p(x) = x - \frac{1}{6}$.
    \item[\ref{pr:fourier}.] The integrals of $\cos(kx)$ over $[-\pi, \pi]$ vanish for $k \neq 0$; $\sin(mx)\cos(nx)$ is an odd function.
\end{description}

\section*{Orthogonal Projections}

\begin{description}[leftmargin=0pt, itemsep=0.3em]
    \item[\ref{pr:proj-line}.] $\pi(\mathbf{x}) = (2, 2, 2)^\top$, $P = \frac{1}{3}\mathbf{1}\mathbf{1}^\top$ (all entries $\frac{1}{3}$), and $\mathbf{x} - \pi(\mathbf{x}) = (0, -1, 1)^\top$.
    \item[\ref{pr:proj-plane}.] (a) $\lambda = \big(\frac{5}{3}, \frac{2}{3}\big)$, $\pi_U(\mathbf{x}) = \frac{1}{3}(7, 2, 5)^\top$, and the distance is $\frac{4}{\sqrt{3}}$. (b) $\mathbf{n} = (-1, 1, 1)^\top$, and $\mathbf{x} - \frac{4}{3} \mathbf{n} = \frac{1}{3}(7, 2, 5)^\top$.
    \item[\ref{pr:centering}.] (a) $C = I_n - P$, where $P = \frac{1}{n}\mathbf{1}\mathbf{1}^\top$ is the projection onto $\Span\{\mathbf{1}\}$. (b) $\frac{1}{n}\mathbf{1}\mathbf{1}^\top \mathbf{x} = \bar{x}\mathbf{1}$; $C\mathbf{x} = (-2, -1, 0, 3)^\top$.
    \item[\ref{pr:distances}.] (a) $1$, closest point $\frac{1}{3}(5, -1, 1)$. (b) $\frac{|10 - 1|}{3} = 3$. (c) The projection is $(2, 2, 2)$, so the distance is $\norm{(-1, 0, 1)} = \sqrt{2}$.
    \item[\ref{pr:least-squares}.] $y = 0.5 + 0.6x$; residuals $(-0.1, 0.3, -0.3, 0.1)$; sum of squared errors $0.2$.
    \item[\ref{pr:ls-special}.] (a) The normal equation is $N c = \sum_i y_i$; the fitted vector $\bar{y}\mathbf{1}$ is the projection of $\mathbf{y}$ onto $\Span\{\mathbf{1}\}$. (b) The normal equation is $(\sum_i x_i^2) w = \sum_i x_i y_i$.
\end{description}

\section*{Rotations}

\begin{description}[leftmargin=0pt, itemsep=0.3em]
    \item[\ref{pr:rot-axis-angle}.] (a) The columns $\frac{1}{3}(2, 2, -1)$, $\frac{1}{3}(-1, 2, 2)$, $\frac{1}{3}(2, -1, 2)$ are orthonormal, and $\det(R) = 1$. (b) $R - I_3$ has the null space spanned by $(1, 1, 1)^\top$, so $\mathbf{v} = \frac{1}{\sqrt{3}}(1, 1, 1)^\top$. (c) $\tr(QR_3(\theta)Q^\top) = \tr(R_3(\theta)Q^\top Q) = \tr R_3(\theta) = 1 + 2\cos\theta$. Here $\tr(R) = 2$, so $\cos\theta = \frac{1}{2}$ and $\theta = 60^\circ$ (up to the sign, which depends on the orientation of the axis).
    \item[\ref{pr:rot-noncomm}.] $R_3(\frac{\pi}{2})R_1(\frac{\pi}{2}) \mathbf{e}_2 = \mathbf{e}_3$ and $R_1(\frac{\pi}{2})R_3(\frac{\pi}{2}) \mathbf{e}_2 = -\mathbf{e}_1$.
    \item[\ref{pr:rot-axis}.] (a) The columns are orthonormal and $\det(C) = 1$. (b) $\mathbf{v} = (1, 1, 1)^\top$. (c) $\tr(C) = 0 = 1 + 2\cos\theta$, so $\theta = 120^\circ$, and three rotations by $120^\circ$ give the identity.
    \item[\ref{pr:givens}.] $\cos\theta = \frac{3}{5}$, $\sin\theta = -\frac{4}{5}$, and $r = 5$.
    \item[\ref{pr:reflections}.] (b) Multiply out and use the addition formulas for $\cos$ and $\sin$.
\end{description}

\section*{The Spectral Theorem}

\begin{description}[leftmargin=0pt, itemsep=0.3em]
    \item[\ref{pr:spectral-2}.] $A_1$: $\lambda = 4, 2$ with $q = \frac{1}{\sqrt{2}}(1, 1)$, $\frac{1}{\sqrt{2}}(1, -1)$; positive definite. $A_2$: $\lambda = 2, -3$ with $q = \frac{1}{\sqrt{5}}(2, 1)$, $\frac{1}{\sqrt{5}}(1, -2)$; not positive definite.
    \item[\ref{pr:spectral-3}.] $A = J + 3I_3$ has the eigenvalues $5$, with $\mathbf{q}_1 = \frac{1}{\sqrt{3}}(1, 1, 1)$, and $2$ (twice), with $\mathbf{q}_2 = \frac{1}{\sqrt{2}}(1, -1, 0)$ and $\mathbf{q}_3 = \frac{1}{\sqrt{6}}(1, 1, -2)$.
    \item[\ref{pr:rayleigh}.] (a) $A = \begin{pmatrix} 0 & \frac{1}{2} \\ \frac{1}{2} & 0 \end{pmatrix}$ has eigenvalues $\pm\frac{1}{2}$, so the maximum is $8 \cdot \frac{1}{2} = 4$ (at $\pm(2, 2)$) and the minimum is $-4$ (at $\pm(2, -2)$). (b) Maximum $3$ at $\pm\frac{1}{\sqrt{2}}(1, 1)$, minimum $-1$ at $\pm\frac{1}{\sqrt{2}}(1, -1)$.
    \item[\ref{pr:quadric}.] The first matrix $\begin{pmatrix} 5 & 2 \\ 2 & 2 \end{pmatrix}$ has eigenvalues $6$ and $1$, with eigenvectors $(2, 1)$ and $(1, -2)$: an ellipse with semi-axes $\frac{1}{\sqrt{6}}$ along $(2, 1)$ and $1$ along $(1, -2)$. The second matrix has eigenvalues $3$ and $-1$: in the rotated coordinates, the curve is $3y_1^2 - y_2^2 = 1$, a hyperbola with axes along $(1, 1)$ and $(1, -1)$.
    \item[\ref{pr:sym-idempotent}.] (b) $P = \sum_{\lambda_i = 1} \mathbf{q}_i \mathbf{q}_i^\top$, and both $\tr(P)$ and $\rank(P)$ equal the number of eigenvalues $1$. (c) $\tr(P) = 1 = \rank(P)$; $P$ projects onto the line $y = x$.
    \item[\ref{pr:laplacian-path}.] (a) $D = \diag(1, 2, 2, 1)$, and $L = \begin{pmatrix} 1 & -1 & 0 & 0 \\ -1 & 2 & -1 & 0 \\ 0 & -1 & 2 & -1 \\ 0 & 0 & -1 & 1 \end{pmatrix}$. (b) $(1 - 2)^2 + (2 - 4)^2 + (4 - 7)^2 = 14$. (c) $Lq = (2 - \sqrt{2})q$, so $q$ is the Fiedler vector; its signs split the path into $\{1, 2\}$ and $\{3, 4\}$.
    \item[\ref{pr:laplacian-cycle}.] (a) Eigenvectors $\frac{1}{2}(1, 1, 1, 1)$ for $0$; $\frac{1}{\sqrt{2}}(1, 0, -1, 0)$ and $\frac{1}{\sqrt{2}}(0, 1, 0, -1)$ for $2$; $\frac{1}{2}(1, -1, 1, -1)$ for $4$. (b) The ratio cuts are $2 \cdot (\frac{1}{2} + \frac{1}{2}) = 2$, $4 \cdot 1 = 4$ and $2 \cdot (1 + \frac{1}{3}) = \frac{8}{3}$, all at least $\lambda_2 = 2$. Equality holds if and only if the vector $f$ of the lecture's theorem on the ratio cut is an eigenvector for $\lambda_2$; here $f = (1, 1, -1, -1)$ for the first split.
    \item[\ref{pr:laplacian-complete}.] (a) $d_i = n - 1$ and $W = \mathbf{1}\mathbf{1}^\top - I_n$. Eigenvalues: $0$ (eigenvector $\mathbf{1}$) and $n$, with multiplicity $n - 1$ (every $x$ with $x^\top\mathbf{1} = 0$). (b) $\mathrm{cut}(A) = |A||\bar{A}|$, so $\mathrm{RatioCut}(A) = |A||\bar{A}|\frac{|A| + |\bar{A}|}{|A||\bar{A}|} = n = \lambda_2$: all splits are equally good, and the graph has no cluster structure.
    \item[\ref{pr:laplacian-components}.] The components are $\{1, 2, 3\}$ and $\{4, 5\}$, so $\Null(L)$ has the basis $(1, 1, 1, 0, 0)$, $(0, 0, 0, 1, 1)$, and $0$ has multiplicity $2$. The block of $\{4, 5\}$ is $\begin{pmatrix} 3 & -3 \\ -3 & 3 \end{pmatrix}$, with the eigenvalue $6$; the block of $\{1, 2, 3\}$ is $\begin{pmatrix} 2 & -2 & 0 \\ -2 & 3 & -1 \\ 0 & -1 & 1 \end{pmatrix}$, with the characteristic polynomial $-\lambda(\lambda^2 - 6\lambda + 6)$, so the remaining eigenvalues are $3 \pm \sqrt{3}$.
\end{description}

\section*{The Singular Value Decomposition}

\begin{description}[leftmargin=0pt, itemsep=0.3em]
    \item[\ref{pr:svd-diag}.] $\sigma_1 = 3$, $\sigma_2 = 2$; for example, $U = \begin{pmatrix} 0 & 1 \\ -1 & 0 \end{pmatrix}$, $\Sigma = \diag(3, 2)$, $V = \begin{pmatrix} 0 & 1 \\ 1 & 0 \end{pmatrix}$.
    \item[\ref{pr:svd-rank1}.] (a) $\sigma_1 = 2$ (and no other nonzero singular value), $\mathbf{u}_1 = \frac{1}{\sqrt{2}}(1, 1, 0)^\top$, $\mathbf{v}_1 = \frac{1}{\sqrt{2}}(1, 1)^\top$. (b) $\col(A)$ is spanned by $\mathbf{u}_1$; $\Null(A^\top)$ by $(1, -1, 0)^\top$ and $\mathbf{e}_3$; $\row(A)$ by $\mathbf{v}_1$; and $\Null(A)$ by $(1, -1)^\top$. (c) $A^+ = \frac{1}{2} \mathbf{v}_1 \mathbf{u}_1^\top = \frac{1}{4}\begin{pmatrix} 1 & 1 & 0 \\ 1 & 1 & 0 \end{pmatrix}$.
    \item[\ref{pr:svd-2x3}.] (a) $AA^\top = \begin{pmatrix} 17 & 8 \\ 8 & 17 \end{pmatrix}$ has eigenvalues $25$ and $9$, so $\sigma_1 = 5$, $\sigma_2 = 3$, $\mathbf{u}_1 = \frac{1}{\sqrt{2}}(1, 1)^\top$, $\mathbf{u}_2 = \frac{1}{\sqrt{2}}(1, -1)^\top$. (b) $\mathbf{v}_1 = \frac{1}{\sqrt{2}}(1, 1, 0)^\top$, $\mathbf{v}_2 = \frac{1}{3\sqrt{2}}(1, -1, 4)^\top$, $\mathbf{v}_3 = \frac{1}{3}(2, -2, -1)^\top$. (c) $A_1 = \begin{pmatrix} 2.5 & 2.5 & 0 \\ 2.5 & 2.5 & 0 \end{pmatrix}$, $\norm{A - A_1}_F = 3$, $\norm{A}_2 = 5$, $\norm{A}_F = \sqrt{34}$.
    \item[\ref{pr:svd-props}.] (a) $A^\top = V\Sigma^\top U^\top$. (b) $A^{-1} = V \Sigma^{-1} U^\top$. (c) $|\det(U)| = |\det(V)| = 1$. (d) $A = \sigma_1 \mathbf{u}_1 \mathbf{v}_1^\top$ has only one singular value; for the example, both norms are $\sqrt{5}$.
    \item[\ref{pr:pca}.] (a) $S = \diag(2, 0.5)$; the principal components are $\mathbf{e}_1$ and $\mathbf{e}_2$; $\frac{2}{2.5} = 80\%$; the average reconstruction error is $0.5 = \lambda_2$. (b) The centered points are $(-2, -2)$, $(0, 0)$, $(2, 2)$, with eigenvalues $\frac{16}{3}$ and $0$: the data lies exactly on a line, and the first principal component explains $100\%$ of the variance. (c) The first right singular vector of the uncentered data points roughly towards the mean $(3, 4)$, not in the direction in which the data varies.
\end{description}

\end{document}
